%-----------PROJECTS-----------
\section{Projects}
    \vspace{2pt}
    \resumeSubHeadingListStart

      \resumeProjectHeading
          {{\textbf{\large{{Few-Shot Bioacoustic Sound Event Detection}}} \space
          \href{https://github.com/nnayz/prototypical-networks-bioacoustics}{\faGithub}}
          $|$ { Torchaudio, PyTorch, PyTorch Lightning, Hydra, MLflow }}{}
          \resumeItemListStart
            \resumeItem{\normalsize{Developed a few-shot meta-learning system to classify animal vocalizations from as few as 5 labeled examples per species, using episodic (N-way, K-shot) training with Prototypical Networks.}}
            \resumeItem{\normalsize{Engineered audio features from raw .wav files into 128-bin log-mel spectrograms and PCEN arrays using librosa, pre-computed and stored as .npy for efficient training.}}
            \resumeItem{\normalsize{Designed and compared four encoder architectures: ResNet baseline (V1), ResNet with channel \& temporal attention + learnable distance metric (V2), Audio Spectrogram Transformer with multi-head self-attention (V3), and pretrained EfficientNet-B1 (V4).}}
            \resumeItem{\normalsize{Built end-to-end training infrastructure with PyTorch Lightning, Hydra configuration management, MLflow experiment tracking, and a custom CLI for training, evaluation, and feature export.}}
          \resumeItemListEnd
          \vspace{2pt}

      \resumeProjectHeading
          {{\textbf{\large{{Federated Fine-Tuning of Nicheformer}}} \space
          \href{https://github.com/nnayz/ft-nicheformer}{\faGithub}}
          $|$ { Python, PyTorch, Federated Learning }}{}
          \resumeItemListStart
            \resumeItem{\normalsize{Implemented federated learning framework for fine-tuning Nicheformer transformer model on mouse brain spatial transcriptomics data, comparing centralized (94.45\%), federated (92.76\%), and local training strategies}}

            \resumeItem{\normalsize{Built data pipeline for anatomical non-IID partitioning (dorsal/mid/ventral regions) with 24-class cell type classification on 250-dimensional features (248 genes + spatial coordinates) achieving competitive 92.76\% accuracy in federated setting}}
          \resumeItemListEnd
          \vspace{2pt}

      \resumeProjectHeading
          {{\textbf{\large{{GitGood: AI-Powered GitHub Repository Analysis Tool}}} \space
          \href{https://github.com/nnayz/gitgood}{\faGithub}}
          $|$ { FastAPI, React, Docker, OpenAI API}}{}
          \resumeItemListStart
            \resumeItem{\normalsize{Built a full-stack web application that integrates with GitHub OAuth to let users import repositories and interact with their codebase through an AI-powered chat interface.}}
            \resumeItem{\normalsize{Designed a FastAPI backend with SQLAlchemy for user session management and repository data storage, connected to the OpenAI API to answer natural language questions about commit history, diffs, and code evolution.}}
            \resumeItem{\normalsize{Developed a responsive React/TypeScript frontend with Vite and TailwindCSS, with Axios handling API communication between client and server.}}
            \resumeItem{\normalsize{Containerised the full stack using Docker and Docker Compose for consistent local development and production deployment.}}
          \resumeItemListEnd
          \vspace{2pt}

      \resumeProjectHeading
          {{\textbf{\large{{ACL Anthology Semantic Retrieval System}}} \space
          \href{https://github.com/nnayz/acl-anthology-rag}{\faGithub}}
          $|$ { FastAPI, Langchain, React, Qdrant, Fireworks AI, Groq }}{}
          \resumeItemListStart
            \resumeItem{\normalsize{Built a Retrieval-Augmented Generation (RAG) system over the ACL Anthology corpus to enable semantic search across tens of thousands of NLP research papers, going beyond keyword matching to understand query meaning using dense vector embeddings.}}
            \resumeItem{\normalsize{Implemented an offline ingestion pipeline that fetches paper metadata and abstracts, generates dense embeddings via Fireworks AI, and indexes them into a Qdrant vector store for fast approximate nearest-neighbour retrieval.}}
            \resumeItem{\normalsize{Designed two query modes: natural language input (reformulated into multiple semantic queries by an LLM via Groq) and "Paper as Query" (using a paper's abstract as a semantic proxy to find related work); results aggregated using Reciprocal Rank Fusion (RRF).}}
            \resumeItem{\normalsize{Developed a FastAPI backend and React frontend with real-time streaming responses, inline citations with hover previews, and a query monitoring panel showing pipeline steps and timing.}}
          \resumeItemListEnd
          \vspace{2pt}

      \resumeProjectHeading
          {{\textbf{\large{{MyTorch: PyTorch-Inspired Deep Learning Framework}}} \space
          \href{https://github.com/nnayz/mytorch}{\faGithub}}
          $|$ { NumPy, Deep Learning}}{}
          \resumeItemListStart
            \resumeItem{\normalsize{Built a minimal deep learning library from scratch inspired by PyTorch's API, implementing core tensor operations (arithmetic, matrix multiplication, transpose) with a custom automatic differentiation engine and backward pass for gradient computation.}}
            \resumeItem{\normalsize{Implemented neural network modules including fully connected (Linear) layers, ReLU activations, and Cross Entropy loss, all with manual forward and backward propagation with no ML framework dependencies.}}
            \resumeItem{\normalsize{Developed SGD and Adam optimisers from first principles, replicating PyTorch's modular optimizer interface.}}
            \resumeItem{\normalsize{Validated the framework on MNIST digit classification, achieving \textbf{97.16\% test accuracy} (macro F1: 97.15\%) using only linear layers, demonstrating correctness of the custom autograd and training pipeline.}}
          \resumeItemListEnd
           \vspace{2pt}

      \resumeProjectHeading
            {{\textbf{\large{{Momentum}}}}
           $|$ { Web Scraping, Content Analysis, Research Tools, Machine Learning, Python}}{}
           \resumeItemListStart
             \resumeItem{\normalsize{Developing a research assistant web application inspired by Anara Labs' Anara, designed to streamline academic and technical research workflows.}}
             \resumeItem{\normalsize{Building automated web content extraction and organization systems to gather, process, and categorize research materials from multiple online sources.}}
             \resumeItem{\normalsize{Implementing intelligent summarization algorithms to condense and present research findings in structured, actionable formats for enhanced productivity.}}
           \resumeItemListEnd

    \resumeSubHeadingListEnd
% \vspace{-12pt}
